\documentclass[UTF8]{ctexart}
\usepackage{amsmath,amssymb}
\usepackage{enumitem}
\usepackage[hidelinks]{hyperref}
\setlist[itemize]{leftmargin=*, itemsep=2pt, topsep=2pt}

\begin{document}

\section*{工作大纲（v9分层修订版：在保持 WS-OVVIS 主线精简优雅的前提下，回挂增强模块与机制储备）}

\subsection*{0. 名称与定位}
本文研究 \textbf{Weakly-supervised Open-Vocabulary Video Instance Segmentation}（WS-OVVIS）的一个严格设定：
\begin{itemize}
\item 训练阶段仅提供 \textbf{clip-level} 正证据标签集合（positive label set），且该集合系统性不完备；
\item 模型必须在不使用像素级实例语义监督、不使用额外图文训练数据的前提下，恢复视频中的实例 mask 序列、track ID 与开放词汇类别；
\item 与经典 OVVIS 相比，本文重点不在实例级强监督下直接学习 ``实例$\leftrightarrow$类别''，而在于：\textbf{给定一个可用的 class-agnostic video instance basis，如何在不完备 clip-level 正证据下进行 open-world 语义归因，并在测试时支持 bag-free 的 text-conditioned open-vocabulary inference。}
\end{itemize}

\par\noindent\hrulefill\par

\subsection*{1. 任务定义与边界说明}

\subsubsection*{1.1 输入、输出与监督边界}
训练输入：
\begin{itemize}
\item 一个原始视频视为一个 clip，记为 $v$；
\item 该 clip 的不完备正证据集合 $Y'(v)$，其中 $Y'(v)\subseteq Y^*(v)$，$Y^*(v)$ 为真实类别集合；
\item 仅允许使用 class-agnostic 伪实例结构监督（默认来自 VideoCutLER 生成的 pseudo mask tubes），不允许使用实例级类别标签、box-level 实例语义监督、像素级类别标注或任何额外图文配对训练数据。
\end{itemize}

测试输入：
\begin{itemize}
\item 视频 clip $v$；
\item 一个文本类别集合 $\mathcal C = \{c_1,\dots,c_K\}$，默认取当前 benchmark 的 base+novel 评测类别全集；
\item \textbf{测试时不输入} $Y'(v)$，也不依赖任何视频内正证据集合做候选约束。
\end{itemize}

输出：
\begin{itemize}
\item clip 内所有预测实例的 mask 序列；
\item clip 内全局 track ID；
\item 每条 track 的开放词汇类别预测；
\item 可选保留 unknown foreground / reject 通道用于诊断，但标准 AP 仍按 top-1 已知类别口径汇报。
\end{itemize}

\subsubsection*{1.2 与经典 OVVIS 的关系}
本文并不改变 OVVIS 的基本输出形式：仍然是 ``视频 + 文本类别空间 $\rightarrow$ 实例分割 + 跟踪 + 类别命名''。不同之处在于：
\begin{itemize}
\item 经典 OVVIS 通常在实例级强监督下直接学习 ``实例$\leftrightarrow$类别'' 关系；
\item 本文仅在训练时获得不完备 clip-level 正证据，因此训练问题必须改写为 \textbf{bag-constrained open-world attribution}；
\item 为尽量对齐经典 OVVIS 的能力边界，测试阶段要求采用 \textbf{bag-free text-conditioned inference}，即不再依赖 $Y'(v)$ 做推理候选约束。
\end{itemize}

\subsubsection*{1.3 open-vocabulary 的主张边界}
\begin{itemize}
\item 主文的核心定义为：\textbf{benchmark-scoped open-vocabulary VIS}，即类别通过文本原型定义，测试时允许识别训练未见的 benchmark novel categories；
\item 主文正式采用 LV-VIS 官方 open-vocabulary 口径：训练期弱监督正证据仅暴露 \textbf{base/seen 641 类}，测试期在 \textbf{base+novel 641+555 类} 全词表上做 bag-free inference；
\item 方法在形式上支持 bag-free 的文本条件推理接口，因此比纯 bag-constrained 分类更接近经典 OVVIS；
\item 但主文不直接主张 unrestricted free-form text grounding，也不把任意外部开放词表识别作为核心结论；若做额外 held-out / out-of-benchmark text 检验，则作为能力边界补充分析，而非主文成立的必要条件。
\end{itemize}

\par\noindent\hrulefill\par

\subsection*{2. 方法总览：压缩后的主线 + 可回挂的机制储备}
整体思想压缩为一条最小闭环：
\begin{itemize}
\item 先利用 class-agnostic 伪实例结构监督训练一个 video-native instance segmentor，获得稳定的 local tracklets；
\item 再在 clip 内把 local tracklets 合并为 global tracks，并用冻结的 DINOv2 特征提取每条 track 的视觉实例表征；
\item 在训练集内部归纳 seen 类的视觉原型，并学习一个强约束的 class-level text-to-prototype map；
\item 训练时在 $Y'(v)+bg+unk$ 上进行 coverage-aware open-world attribution，避免 hidden positives 被压成背景；
\item 测试时丢弃 bag 约束，直接用 ``track DINO feature $\leftrightarrow$ mapped text prototype'' 的全词表匹配进行 bag-free open-vocabulary inference。
\end{itemize}

\subsubsection*{2.1 class-agnostic video instance basis}
\begin{itemize}
\item 主线固定使用 \textbf{VideoCutLER} 生成 class-agnostic pseudo mask tubes，作为结构监督来源；所有 baseline、主方法与消融共享同一份 pseudo tubes，以保证公平性。
\item 伪 tube 的预处理仅保留与结构质量直接相关的步骤：短轨过滤、面积稳定性过滤、时序一致性过滤、重叠实例去重以及低置信实例剔除；最终统一导出每帧 mask、tube ID 与 tube confidence。
\item 采用 SeqFormer 一类 video-native instance segmentor 学习 local instance masks、tracklets 与跨帧实例 query；训练目标仅包含 class-agnostic mask / tracking / association，不使用实例级类别监督。
\item 该阶段的唯一目标是获得时序更稳定、mask 更平滑、碎片更少的实例 basis；此处 \textbf{不承担开放词汇分类功能}。
\item 对后续模块的标准导出单元固定为 local tracklet $k$：
\[
(q_k^{loc}, M_k^{loc}, s_k^{mask}, t_k^{start}, t_k^{end}),
\]
其中分别表示局部 query、局部 mask 序列、平均 mask score 与时间起止位置。后续 2.2--2.6 统一消费这一接口。
\end{itemize}

\subsubsection*{2.2 clip-level global track bank}
\begin{itemize}
\item 将一个原始视频视为一个 clip；前向计算时按固定长度的重叠 window 切分，默认使用 \textbf{$T=24$ 帧、stride $=12$ 帧}。window 仅用于计算，不携带独立语义标签。
\item 在相邻重叠 window 之间，把 local tracklets stitching 成 clip-level global tracks，得到 $\mathcal T(v)=\{\tau\}_{\tau=1}^{N_v}$；后续所有语义建模、归因与评测统一以 clip-level global track 为对象。
\item stitching 主线保持 class-agnostic，仅使用时空几何重叠与局部 query 一致性，不让语义分数反向主导 track 结构。对相邻 window 中 local tracks $k,\ell$，匹配分数定义为
\[
 s(k,\ell)=\lambda_1 \overline{\mathrm{IoU}}_{overlap}+\lambda_2 \cos(q_k^{loc},q_\ell^{loc}).
\]
\item 在得分超过阈值的候选 pair 上做 Hungarian matching，构造 global tracks。主线默认离线固定一份 clip-level track bank，用于归因训练与测试，避免训练期同时漂移 ``数据平面'' 与 ``归因目标''。
\item 对长视频若局部 linking 误差明显累积，可在附录中增加两层式 segment-to-clip stitching；但主文不采用全局图优化，也不在 linker 中额外加入语义 rerank。
\end{itemize}

\subsubsection*{2.3 DINOv2 track feature：唯一语义主载体}
\begin{itemize}
\item 主线只强调一条原则：\textbf{track-level DINO 特征 $z_\tau$ 是后续开放词汇语义建模的唯一主语义载体}。query 聚合得到的 $q_\tau$ 仅保留为 linking 一致性与误差诊断的辅助量，不进入主语义桥。与 $q_\tau$ 相关的预归因启发式和诊断信号统一下放到 2.7.1，作为开发期可选储备。
\item DINO 主干固定为 \textbf{DINOv2 ViT-B/14}。对每条 global track，在其可见帧上根据实例 bbox 做裁剪，并在 bbox 周围保留约 \textbf{10\% 上下文 padding}；随后对 crop 内 patch 特征做 \textbf{mask-aware patch pooling}，得到逐帧实例特征 $\{d_\tau^t\}$。
\item 在可见帧上按 mask 质量与可见性加权聚合，得到 track-level 视觉特征：
\[
 z_\tau = \frac{\sum_{t\in \mathcal V(\tau)} w_{\tau,t}^{d}\, d_\tau^t}{\sum_{t\in \mathcal V(\tau)} w_{\tau,t}^{d}},
\]
其中 $w_{\tau,t}^{d}$ 由该帧 mask score、可见面积与时序稳定性联合给出。主线默认采用 visible-frame weighted mean pooling，不引入 attention 或复杂 memory 聚合。
\item 同时为每条 track 定义前景质量分数 $o_\tau^{obj}$，仅依赖结构侧统计量：
\[
 o_\tau^{obj} = \sigma\!\left(\gamma_1 \, \overline{s_\tau^{mask}} + \gamma_2 \,\mathrm{dur}_\tau + \gamma_3 \,\mathrm{cons}_\tau\right),
\]
其中分别对应跨帧平均 mask score、可见帧占比与时序一致性。该量仅用于区分 ``可解释前景'' 与 ``低质量伪结构''，不直接承担类别判别。
\end{itemize}

\subsubsection*{2.4 训练集内部类视觉原型与 class-level text map}
\begin{itemize}
\item 文本侧固定使用 \textbf{冻结 CLIP Text Encoder + 类名 only}。主文默认不启用 prompt ensemble、同义词扩展或 LLM 描述增强；这些仅作为附录分析。
\item 不使用任何额外图文数据，也不在 noisy instance 上直接学习自由的 text$\rightarrow$DINO 大桥。相反，先在训练集内部归纳 \textbf{seen/base 类视觉原型}。对每个训练 seen 类 $y$，在所有满足 $y\in Y'(v)$ 的 clip 中，依据高 objectness 与 clip-level 正证据覆盖启发式选择 latent positive tracks，并据此得到 DINO 空间中的视觉原型 $p_y$。
\item prototype bank 的主线只保留两步：\textbf{先离线初始化一版 prototype bank；训练中按低频策略刷新。} 不把 EM 轮数、刷新间隔、momentum 细节放进核心主张；这些实现细节在 2.7.2 中以 ``prototype 机制储备 的形式保留，便于开发时快速回挂。
\item 在此基础上，仅学习一个 \textbf{强约束的 class-level map} $A$，使文本嵌入 $t_y$ 能映射到视觉原型附近：
\[
\hat p_y = \mathrm{LN}(A\,\mathrm{LN}(t_y)).
\]
主实现采用 \textbf{单层线性映射 $A$ + 强 L2 正则}，并使用正交 Procrustes 作为初始化；不使用大 MLP，不做 instance-level 自由文本桥接。
\item 对于训练中有视觉原型的 base 类，$\hat p_y$ 由 $At_y$ 与 $p_y$ 的一致性共同约束；对于 novel 类，则直接使用 $\hat p_j=\mathrm{LN}(A\,\mathrm{LN}(t_j))$ 作为 DINO 空间中的文本预测原型。
\end{itemize}

\subsubsection*{2.5 训练期：coverage-aware open-world attribution（核心主线）}
\textbf{输入：}
\begin{itemize}
\item 当前 clip 的 DINO track feature set $Z(v)=\{z_\tau\}$ 与对应前景质量分数 $\{o_\tau^{obj}\}$；
\item 当前 clip 的正证据集合 $Y'(v)$；
\item mapped text prototypes $\hat P=\{\hat p_j\}$。
\end{itemize}

\textbf{主线建模原则：}
\begin{itemize}
\item 训练只在 \textbf{$Y'(v)+bg+unk$} 上进行归因：$Y'(v)$ 用于承接已观测到的正证据，$bg$ 用于吸收低质量或非前景实例，$unk$ 用于承接 ``当前未被正证据解释、但又不应塌成背景'' 的前景质量；
\item 核心不是让每条 track 必然落到某个已知类，而是让 ``已知正证据得到覆盖'' 与 ``hidden positives 不被压成背景'' 同时成立；
\item 因此，主方法必须显式保留 \textbf{unknown foreground slack} 与 \textbf{coverage constraint}。
\end{itemize}

\textbf{相似度与代价：}
\begin{itemize}
\item 对 $j\in Y'(v)$，定义语义相似度与 logit：
\[
 s_{\tau j}=\cos(z_\tau,\hat p_j),\qquad l_{\tau j}=s_{\tau j}/T_{temp}.
\]
\item 对语义列 $j\in Y'(v)$，定义代价 $C_{\tau j}=-l_{\tau j}$。
\item 背景列与 unknown 列的代价仅由结构质量控制：
\[
 C_{\tau,bg}=c_{bg}-\lambda_{bg}(1-o_\tau^{obj}),
\]
\[
 C_{\tau,unk}=c_{unk}-\lambda_{unk}o_\tau^{obj}.
\]
其含义是：低 objectness 的 track 更容易被解释为 bg，而高 objectness 但无法由当前正证据解释的 track 允许流向 unk。
\end{itemize}

\textbf{归因主方法：}
\begin{itemize}
\item 主实现采用 coverage-aware entropic soft assignment，在 $Y'(v)\cup\{bg,unk\}$ 上求软分配矩阵 $P^{(v)}$：
\[
 \min_{P^{(v)}\ge 0}\ \langle C,P^{(v)}\rangle - \varepsilon H(P^{(v)})
 + \lambda_{\mathrm{row}}\,\mathrm{KL}(P^{(v)}\mathbf 1\|a^{(v)})
 + \lambda_{\mathrm{cov}}\sum_{y\in Y'(v)} \Big[m_y^{(v)}-\sum_\tau P^{(v)}_{\tau y}\Big]_+^2.
\]
\item 其中 $a_\tau^{(v)}\propto o_\tau^{obj}$ 并归一化，表示每个 track 需被解释的质量；coverage 下界采用统一低下界
\[
 m_y^{(v)}=\rho_{cov}\cdot \frac{1}{N_v}\sum_\tau a_\tau^{(v)},\qquad y\in Y'(v),
\]
仅用于防止正证据被完全漏掉，而不编码每类实例数先验。
\item 该形式保留了 many-to-one、多实例共享同类、open-world slack 与正证据覆盖的核心思想，同时避免把训练主线建成过度复杂的多重候选检索系统。
\item 求解上采用熵正则的可微迭代更新；主线默认仅对 class-level map $A$、prototype 更新规则与归因相关标度参数进行优化，不端到端更新整个 SeqFormer 主干与 DINOv2 特征提取器。
\end{itemize}

\textbf{训练注入：}
\begin{itemize}
\item 仅当某 track 在非 bg/unk 列上的总质量 $s_\tau^{sem}=\sum_{j\in Y'(v)}P_{\tau j}^{(v)}$ 超过阈值 $\tau_{sem}$ 时，才将 $P^{(v)}$ 在非 bg/unk 列上按行归一化为软标签，并以 $\mathrm{softmax}(l_{\tau j})$ 与该软标签做 KL/CE 监督；否则该 track 不施加语义 CE。
\item 训练启用顺序固定为：\textbf{closed-world bag baseline $\rightarrow$ open-world（$Y'(v)+bg+unk$）$\rightarrow$ +coverage penalty}。retrieval、warm-up BCE 与额外一致性正则全部降为 2.7 中的增强模块，不进入主线成立条件。
\end{itemize}

\subsubsection*{2.6 测试期：bag-free text-conditioned open-vocabulary inference}
\begin{itemize}
\item 测试阶段不再使用 $Y'(v)$，也不再围绕视频内正证据构造候选集。给定视频与文本类别集合 $\mathcal C$，对每条 global track 与全部文本类别直接计算
\[
 s_{\tau,c}=\cos(z_\tau,\hat p_c),\qquad c\in\mathcal C.
\]
\item 推理采用与训练一致的 sliding-window + clip-level stitching 流程，仍使用 2.2 中相同的 class-agnostic linker，避免推理阶段引入额外语义 rerank 破坏结构口径。
\item 主结果默认输出 top-1 known class；unknown foreground 分支主要用于诊断高 objectness 但低文本可解释性的 tracks，不作为标准 AP 的强依赖后处理。
\item 若需要进一步分析 unknown 通道的作用，可在附录中汇报 ``被 unk 拦截的比例''、``强制回填为 best-known 时的 AP 变化'' 以及 ``unk 对 HPR / UAR 的影响''；但主文不把复杂的 unk-fallback 逻辑作为核心方法的一部分。
\end{itemize}

\subsubsection*{2.7 可选增强模块与机制储备（开发备用层）}
2.1--2.6 构成论文主线；本节则把从 v7 中移出的防御性机制、实现细节与备用方案系统化收纳为 \textbf{开发备用层}。它们的存在目的不是让主文重新变厚，而是：当沿核心主线开发时遇到具体失败模式，可以按触发条件快速启用相应机制，而不必重新发明备选方案。

\subsubsection*{2.7.1 track-side 辅助信号储备：$q_\tau$、$r_\tau^{pre}$ 与全词表诊断}
\begin{itemize}
\item 主线不依赖 $q_\tau$ 做语义桥，但开发中仍建议保留一套轻量 track-side 辅助信号。对每条 global track，可继续定义 query 聚合：
\[
 q_\tau = \frac{\sum_{k\in \mathcal S(\tau)} w_k^{q} \, q_k^{loc}}{\sum_{k\in \mathcal S(\tau)} w_k^{q}},
\]
其中 $w_k^{q}$ 由 tracklet 长度、平均 mask score 与 overlap 关联质量联合给出。该量默认只服务于 linking 一致性、错误诊断与可选 rerank，不进入主语义分类头。
\item 若开发中发现 unk/bg 分流过于依赖单一 objectness，可恢复一组全词表预归因启发式。设 mapped text prototypes 为 $\hat P=\{\hat p_j\}$，则全词表开放词汇 logit 为
\[
 l_{\tau j}^{full}=\cos(z_\tau,\hat p_j)/T_{temp}.
\]
在 full-vocab softmax 上取 top-1 概率 $p_\tau^{full}$ 与 top1-top2 margin $m_\tau^{full}$，定义
\[
 r_\tau^{pre}=\sigma\!\left(\xi_1 p_\tau^{full}+\xi_2 m_\tau^{full}\right).
\]
该定义不依赖候选集合，可作为 ``当前 track 是否已被已知类别较好解释'' 的先验诊断量。
\item 推荐用途：\textbf{(i)} 诊断 ``高 objectness 但低文本可解释性'' 的轨迹；\textbf{(ii)} 在 2.7.3 的 retrieval 中做高置信前景筛选；\textbf{(iii)} 在 2.7.5 的 unk-fallback 中辅助判定是否回填为 best-known。若核心版训练已稳定，则不必默认启用这些附加信号。
\end{itemize}

\subsubsection*{2.7.2 prototype 机制储备：latent positive 挖掘、EM 与低频刷新}
\begin{itemize}
\item 主线只要求 ``prototype bank 可初始化且可低频刷新''；开发备用层则保留更具体的原型构建程序。对每个训练 seen 类 $y$，在所有满足 $y\in Y'(v)$ 的 clip 中，用高 objectness、高时序稳定与 bag-level coverage 启发式选出初始 latent positive tracks，再通过 EM 风格的 ``选 latent positive tracks $\rightarrow$ 更新类视觉原型'' 过程得到 DINO 空间中的视觉原型 $p_y$。
\item 推荐的保守实现为：\textbf{初始原型做 3 轮 EM}；进入归因训练后，\textbf{每 2 个 epoch 刷新一次 prototype}，并采用 \textbf{momentum update} 平滑更新，以避免 prototype 与 attribution 同时剧烈漂移。
\item class-level map $A$ 的主实现仍采用 \textbf{单层线性映射 + 强 L2 正则 + 正交 Procrustes 初始化}；ridge-style 闭式映射和纯 Procrustes 保留为等预算替代实现，用于判断映射是否被训练噪声带偏。
\item 触发条件：若核心版出现 ``prototype 漂移快、base 类语义桥不稳、训练早期 AP/HPR 波动大''，优先从这一储备层恢复更保守的 EM + 低频刷新 + momentum 组合，而不是先增加更复杂的语义损失。
\end{itemize}

\subsubsection*{2.7.3 candidate retrieval 储备：缓解 hidden class 漏候选}
\begin{itemize}
\item 若核心版在 high-objectness tracks 上频繁出现 ``明明像前景，但 $Y'(v)$ 中没有可解释类，导致大量质量被挤向 unk''，可启用 clip-level retrieved candidates 作为增强。先计算开放词汇相似度
\[
 s_{\tau j}=\cos(z_\tau,\hat p_j),\qquad l_{\tau j}=s_{\tau j}/T_{temp}.
\]
\item 仅对高前景置信轨迹参与检索，设 $\mathcal I_{fg}(v)=\{\tau\mid o_\tau^{obj}>\tau_{fg}\}$；对每个类别 $j$ 定义 clip 级检索分数
\[
 u_j(v)=\operatorname{LSE}_{\tau\in \mathrm{Top}\text{-}r_j(v),\ \tau\in\mathcal I_{fg}(v)}(o_\tau^{obj}\, l_{\tau j}),
\]
其中主实现默认 \textbf{$r=3$}。候选列集合定义为
\[
 \hat Y(v)=Y'(v)\cup \operatorname{TopK}_{j\notin Y'(v)}(u_j(v)),
\]
默认 \textbf{$K_{\max}=5$}。
\item 推荐采用 \textbf{递进式启用}：训练初期只用 $Y'(v)$；待 prototype 与 $A$ 稳定后，再打开 retrieved candidates。这样可以把 ``open-world 核心收益'' 与 ``候选扩展收益'' 分开看，避免过早把主线重新变成复杂候选系统。
\item 主文写作原则不变：先报告不含 retrieval 的 core，再报告 ``core + retrieval''。开发上则把 retrieval 视为 \textbf{hidden positives 漏候选问题} 的第一备用解，而非默认必选项。
\end{itemize}

\subsubsection*{2.7.4 训练稳定器储备：warm-up BCE、一致性与轻量 joint fine-tune}
\begin{itemize}
\item \textbf{bag-level warm-up BCE}：若 open-world transport 在训练最初期不稳定，可在前 10\%--20\% 步数内，对 $Y'(v)$ 上类别使用 clip-level top-r LSE 聚合的 BCE 作为 warm-up；中后期关闭，避免与 transport 目标重复并重新压制 hidden positives。
\item \textbf{时序语义一致性}：若同一 clip-level track 在不同增强视图或相邻计算片段中的语义分布波动过大，可加 KL 或 InfoNCE 式的一致性弱正则；前提是 2.2 的 stitching 已基本稳定，否则该项可能把 linking 噪声固化。
\item \textbf{轻量 joint fine-tune}：主线默认冻结 backbone、pixel decoder、大部分 mask head 与 DINOv2 特征提取器；若后期确认瓶颈来自表示错位而非伪结构噪声，可只对极小范围参数做轻量 joint fine-tune，但始终应把它视为附录增强，不与 core 主张混为一谈。
\item 启用优先级建议：先尝试 warm-up BCE，再尝试一致性正则，最后才考虑轻量 joint fine-tune；不要在核心版尚未站稳时同时打开多种稳定器。
\end{itemize}

\subsubsection*{2.7.5 推理期储备：unknown fallback 与诊断性拦截}
\begin{itemize}
\item 主结果默认输出 top-1 known class，但开发中仍可保留一套更完整的 unk 处理逻辑。已知类、unknown 与背景的判决可结合 $o_\tau^{obj}$、top-1 置信度与 top1-top2 margin 完成；若某条 global track 在测试时 unk 得分最高，则仅当其 best-known class 置信度超过最小阈值 $\tau_{known}$ 时，才回填为该 known class；否则仅保留在诊断分支中，不计入标准 AP 输出。
\item 建议同时汇报：\textbf{(i)} 被 unk 拦截的 track 比例；\textbf{(ii)} 回填比例；\textbf{(iii)} 将这些 track 强制回填为 best-known 时的 AP 变化。这样可以更清楚地区分 ``unk 真的保护了 hidden positives'' 与 ``unk 只是过度保守地吃掉已知类''。
\item 启用条件：当 HPR 提升但 AP 明显下降，且定性分析显示大量高质量轨迹停留在 unk 分支时，再打开 fallback 逻辑做阈值分析；若 core 已能平衡 AP 与 HPR，则不必在主结果中引入复杂回填策略。
\end{itemize}

\subsubsection*{2.7.6 one-round quality-aware refinement 储备}
\begin{itemize}
\item 在获得第一轮 global tracks 与语义归因后，可用 \textbf{one-round quality-aware refinement} 对伪实例结构进行纠错；默认只做 \textbf{1 轮 refinement}，且默认 \textbf{不启用 label-set expansion，不重建 clip-level track bank}，以免方法重新演化成多轮自训练系统。
\item refinement 的最小更新单元是 clip-level global track，而不是 window 内 local track。对每条 global track 采用分解式质量评估：
\[
 Q_\tau^{keep}=w_m Q_\tau^{mask}+w_t Q_\tau^{temp}+w_a Q_\tau^{agree},\qquad
 Q_\tau^{diag}=Q_\tau^{sem}.
\]
其中 $Q_\tau^{mask}$ 由跨帧平均 mask score 构成，$Q_\tau^{temp}$ 由相邻帧及重叠片段上的时序一致性构成，$Q_\tau^{agree}$ 由与旧伪 tube 的 IoU/关联一致性构成，$Q_\tau^{sem}$ 由 top-1 概率与 margin 构成。
\item keep/drop 的主依据仍然是 mask 质量与时序稳定性，而不是语义分数。第一轮 refinement 中，global track 的 keep/drop \textbf{完全由 $Q_\tau^{keep}$ 决定}；$Q_\tau^{diag}$ 只用于命名辅助和 unknown 诊断，避免噪声文本桥反向破坏 class-agnostic instance basis。
\item 伪标签更新规则采用 decoupled gating：
\begin{itemize}
\item keep：若 $Q_\tau^{mask}>\tau_m$ 且 $Q_\tau^{temp}>\tau_t$，并且与旧伪 tube 的一致性不低于阈值，则保留为下一轮伪 track；
\item drop：若 $Q_\tau^{mask}<\tau_m^{-}$ 或 $Q_\tau^{temp}<\tau_t^{-}$，则直接丢弃；
\item semantic-uncertain keep：对 mask/temporal 质量高但语义不稳定的 track，保留为 class-agnostic 伪 track，不因语义弱而删除。
\end{itemize}
\item mask 更新采用旧 pseudo 与当前预测的 logit 级加权融合：
\[
\begin{aligned}
 \tilde M_\tau &= \lambda_{old} M_\tau^{old}+\lambda_{new} M_\tau^{pred},\\
 \lambda_{new} &= \mathrm{clip}\!\left(
 \frac{Q_\tau^{mask}+Q_\tau^{temp}}{Q_\tau^{mask}+Q_\tau^{temp}+Q_\tau^{agree}+\epsilon}, \lambda_{\min},\lambda_{\max}
 \right).
\end{aligned}
\]
其中推荐默认 $\lambda_{\min}=0.2,\lambda_{\max}=0.8$。refinement 后的训练从上一阶段 checkpoint 继续短 schedule 微调，而不是从头重训；语义桥参数继承上一轮，但学习率单独设小，以减少伪标签刷新带来的分布震荡。
\item 启用条件：当核心版的主要瓶颈来自 \textbf{mask 破碎、时序粘连、track 结构脏}，而不是开放词汇归因本身时，再打开这一模块。也就是说，refinement 是针对 ``结构层失败'' 的备用方案，不是主线语义问题的第一修复手段。
\end{itemize}

\subsubsection*{2.7.7 开发启用顺序：从核心闭环到备用层的升级路径}
\begin{itemize}
\item \textbf{Level 0 / Core only}：先只跑 2.1--2.6，确认最小闭环已经可训练、可推理，并得到可解释的 AP / HPR / UAR 基线。
\item \textbf{Level 1 / 语义桥保守化}：若训练早期波动大，优先回挂 2.7.2 的 EM 初始化、低频 prototype 刷新与 momentum 更新；必要时加 2.7.4 的 warm-up BCE。
\item \textbf{Level 2 / hidden positives 承接增强}：若高 objectness 轨迹大量流向 unk，且定性上存在明显漏候选，再启用 2.7.3 的 retrieved candidates；若 unk 过度保守，再配合 2.7.5 做 fallback 阈值分析。
\item \textbf{Level 3 / 结构层修补}：若语义桥已基本成立，但 track 仍明显破碎或黏连，再启用 2.7.6 的 one-round refinement。
\item \textbf{Level 4 / 仅在必要时追加稳定器}：时序一致性、轻量 joint fine-tune 等放在最后；除非已经明确定位到相应故障源，否则不要同时开启多个防御机制。
\end{itemize}

\par\noindent\hrulefill\par

\subsection*{3. 论文核心主张与贡献（修订后）}
\begin{itemize}
\item 提出一个严格的 WS-OVVIS 设定：训练仅给不完备 clip-level 正证据集合，要求恢复视频中的实例轨迹与开放词汇类别；
\item 将训练问题显式写成 \textbf{coverage-aware bag-constrained open-world attribution}，把 ``hidden positives 不应塌成背景'' 作为核心科学问题，而非附属现象；
\item 提出一条更凝练的语义主线：\textbf{DINO track feature + visual class prototype + class-level text map}，在不使用额外图文训练数据的前提下实现开放词汇语义桥；
\item 明确区分训练期的 bag-constrained learning 与测试期的 bag-free text-conditioned inference，使方法在更弱训练监督下尽量对齐经典 OVVIS 的推理接口；
\item 在方法组织上明确区分 \textbf{核心闭环}、\textbf{增强模块} 与 \textbf{机制储备}：主文保持精简，而开发文档仍保留防御性方案与实现细节，方便沿主线失败时快速回挂备用机制。
\end{itemize}

\par\noindent\hrulefill\par

\subsection*{4. 实验设计}

\subsubsection*{4.1 数据集与任务口径}
\begin{itemize}
\item 主数据集：\textbf{LV-VIS}，作为 open-vocabulary 主战场。正式采用官方划分：train/val/test = 3083/837/980 videos，base/novel = 641/555 categories。主结果报告 overall mAP、$AP_{base}$、$AP_{novel}$。
\item clip 定义：\textbf{一个原始视频 = 一个 clip}。训练与推理的窗口切分仅用于计算，固定采用 $T=24$、stride $=12$。
\item 弱监督正证据构造：对训练集每个 clip，由其真实类别集合 $Y^*(v)$ 构造不完备集合 $Y'(v)$；训练期仅使用 $Y'(v)$，val/test 评测仍使用原始 GT，不做缺失。
\item 主文保留两种缺失协议：\textbf{Uniform Missing} 与 \textbf{Long-tail Missing}；\par\hspace*{1.6em}\textbf{Scenario/Domain Missing} 仅放附录。
\item missing rate 扫描点固定为 \textbf{$p\in\{0,0.2,0.4,0.6\}$}；主表默认使用 \textbf{$p=0.4$}，并同时报告完整 robustness curve 与 AURC。
\item 训练期仅在 LV-VIS train 上构造弱监督正证据集合 $Y'(v)$；测试时统一采用 bag-free inference，不输入 $Y'(v)$。
\item 泛化评测：在 LV-VIS train 上完成训练后，直接 zero-shot transfer 到 \textbf{YouTube-VIS 2019/2021、OVIS、BURST}；主文仅保留摘要性泛化结果，详细表格放附录。
\item 若做 out-of-benchmark text 检验，则只作为边界分析，不与主文的 benchmark scoped 结论混为一谈。
\end{itemize}

\subsubsection*{4.2 评测指标}
\begin{itemize}
\item 标准 OVVIS 指标：LV-VIS 上的 mAP / $AP_{base}$ / $AP_{novel}$；YTVIS/OVIS 上的 mAP；BURST 上的 mAP 与 HOTA。
\item 弱监督不完备标签诊断：
\begin{itemize}
\item \textbf{Set-Coverage Recall (SCR)}：对每个 clip、每个 $y\in Y'(v)$，若存在至少一条预测 global track 与某个 GT 中类别为 $y$ 的实例匹配且 stIoU $\ge 0.5$，则认为该正证据被覆盖；再对 $Y'(v)$ 求平均。
\item \textbf{Hidden-Positive Recall (HPR)}：对 $Y^*(v)\setminus Y'(v)$ 对应实例，统计其是否被任意非 bg 预测 global track 承接，用于衡量 hidden positives 是否未被压成背景。
\item \textbf{Unknown Assignment Recall (UAR)}：在被承接的 hidden positives 中，统计被 unk 通道承接的比例，用于验证 unknown foreground slack 的必要性。
\item \textbf{Missing-label robustness curve}：报告 mAP 随 missing rate $p$ 的变化曲线，并计算 \textbf{AURC}（Area Under Robustness Curve）。
\end{itemize}
\item 论文成立的关键证据是：在标准 AP 不显著退化的前提下，HPR / UAR / AURC 相对 closed-world weak baseline 有稳定提升。
\end{itemize}

\subsubsection*{4.3 训练流程与公平性口径}
\begin{itemize}
\item 所有 baseline、主方法与消融共享同一份 VideoCutLER pseudo tubes、相同训练视频、相同 clip-level 标签协议、相同文本词表与相同 DINO/CLIP Text 先验预算。
\item 训练分段为：
\begin{itemize}
\item Stage A：生成并过滤 pseudo mask tubes；
\item Stage B：用 pseudo tubes 训练 class-agnostic SeqFormer，导出 local tracklets；
\item Stage C：构建 clip-level track bank，提取 DINO track feature，初始化 prototype bank，并训练核心版 open-world attribution 与 class-level map；
\item Stage D（可选增强）：按 2.7.7 的启用顺序选择性加入 prototype 保守刷新、retrieval、unk-fallback 分析、refinement 与其他稳定器，并做短 schedule 微调。
\end{itemize}
\item 主线默认冻结 backbone、pixel decoder、大部分 mask head 与 DINOv2 特征提取器，仅优化 class-level map $A$、prototype 更新规则与归因相关标度参数；可选的轻量 joint fine-tune 仅作为附录增强。
\item 开发建议始终遵循 ``先 core、后储备 的顺序：任何增强模块都应以 2.1--2.6 的最小闭环已经可复现为前提，再根据 2.7.7 的故障定位结果逐项启用。
\item 与经典 OVVIS 直接比较绝对 AP 不是严格公平的，因为其训练通常使用更强实例级监督；因此经典 OVVIS 方法只能作为 \textbf{reference / upper bound}，不能充当主表的同预算公平对手。
\item 真正公平的比较应满足：相同训练视频、相同 clip-level 标签协议、相同外部先验预算、且不允许额外像素级实例语义监督。本文与经典 OVVIS 的公平性主张应表述为：\textbf{训练监督更弱，但测试接口与能力边界尽量对齐。}
\end{itemize}

\subsubsection*{4.4 必须比较的 baseline}
实验建议分成三层：
\begin{itemize}
\item \textbf{主表（公平比较）}：只放与本文训练信息量尽量一致的方法；
\item \textbf{主线拆解表}：专门比较核心版与增强版，证明论文成立不依赖额外稳定器；
\item \textbf{参考表（能力边界 / 强监督上界）}：放经典 OVVIS 方法，作为参考而非主公平对手。
\end{itemize}

主表必须包含：
\begin{itemize}
\item raw pseudo-tube / raw proposal + text matching lower bound；
\item SeqFormer-on-pseudo-tubes，无 prototype、无 open-world attribution 的 class-agnostic baseline；
\item closed-world weak baseline（只允许 $Y'(v)+bg$，无 unknown）；
\item query-only 语义桥 baseline 与 DINO track feature 主线的对照；
\item no-prototype / no-map / no-bag-free-inference 等关键消融；
\item 至少一个弱监督 VIS 代表方法的适配版（如 MaskConsist-style baseline）。
\end{itemize}

主线拆解表至少包含：
\begin{itemize}
\item core：2.1--2.6 的最小闭环版本；
\item core + retrieval；
\item core + one-round refinement；
\item core + retrieval + refinement（完整增强版）。
\end{itemize}

参考表建议包含：
\begin{itemize}
\item OV2Seg；
\item OpenVIS / InstFormer；
\item OVFormer；
\item CLIP-VIS；
\item BriVIS；
\item 若涉及实时性，再补 TROY-VIS。
\end{itemize}

\subsubsection*{4.5 关键消融与能力边界分析}
\begin{itemize}
\item DINO track feature vs query-only 语义主载体；
\item 有无 visual class prototypes；
\item 有无 class-level map $A$；
\item class-level map 的实现形式：线性层 + L2 正则 vs 纯 Procrustes vs ridge-style 闭式映射；
\item core attribution 的各组成部分：unknown、coverage；
\item retrieved candidates 是否值得进入完整系统：无 retrieval vs +TopK retrieval；
\item testing bag-free inference vs testing 仍受 bag 约束；
\item held-out text / synonyms / prompt ensemble 对 bag-free inference 的影响；
\item no refinement vs one-round refinement；
\item prototype 刷新策略、warm-up BCE、时序一致性、unk fallback 阈值与 refinement 细节等全部作为第二层消融，用于验证机制储备的启用条件，而不与核心命题混在同一优先级。
\end{itemize}

\subsubsection*{4.6 定性与诊断分析}
\begin{itemize}
\item 归因矩阵 $P^{(v)}$ 的可视化（track $\times$ label），展示从 closed-world 到 open-world、再到 +coverage 的变化；
\item bg/unk 分配与 transport mass 的可视化，验证正证据覆盖与集合外前景承接是否同时成立；
\item core 版本与完整增强版在同一视频上的对照案例，展示哪些收益来自核心方法、哪些收益来自增强模块；
\item 失败模式分析：遮挡交换、同类多实例混淆、小目标与长尾类别；
\item Missing-label 失败模式：当 $Y'(v)$ 丢失关键类时，模型是否错误压 bg、是否由 unk 承接、coverage 是否误触发导致假阳性。
\end{itemize}

\par\noindent\hrulefill\par

\subsection*{5. 预期结论组织方式}
\begin{itemize}
\item 第一层结论：在更弱监督、更不完备标签的设定下，本文的 \textbf{核心闭环} 已能建立可用的 open-vocabulary video instance basis，并显著缓解 hidden positives 塌向背景的问题；
\item 第二层结论：测试期的 bag-free inference 使方法的接口与能力边界明显接近经典 OVVIS，而不仅仅是一个视频内 bag 归因器；
\item 第三层结论：retrieval、refinement 与其他稳定器能够带来进一步收益，但这些收益建立在核心闭环已成立的基础上，而不是替代核心命题；机制储备的价值在于让开发过程具备可控的升级路径，而不是让主文重新回到复杂堆叠。
\item 因而，本文的核心价值不在于 ``完全替代强监督 OVVIS''，而在于：\textbf{在不完备 clip-level 正证据下，以更弱训练监督逼近经典 OVVIS 的推理能力边界，并以更清晰、更简洁的技术主线支撑这一主张；而分层保留的机制储备则保证了实际开发中遇到困难时仍有系统化备用方案可调用。}
\end{itemize}

\par\noindent\hrulefill\par

\end{document}
